% Generated draft: requires mathematical and source review.
\documentclass[11pt]{article}
\usepackage[margin=25mm]{geometry}
\usepackage{fontspec}
\setmainfont{Latin Modern Roman}
\usepackage{amsmath,amssymb,mathtools,unicode-math}
\setmathfont{Latin Modern Math}
\usepackage{enumitem,needspace,xurl,hyperref}
\setlength{\parindent}{0pt}
\setlength{\parskip}{5pt}
\setlength{\emergencystretch}{3em}
\newcommand{\N}{\mathbb{N}}
\newcommand{\Z}{\mathbb{Z}}
\newcommand{\Q}{\mathbb{Q}}
\newcommand{\R}{\mathbb{R}}
\newcommand{\C}{\mathbb{C}}
\newcommand{\F}{\mathbb{F}}
\newcommand{\A}{\mathbb{A}}
\newcommand{\PP}{\mathbb{P}}
\newcommand{\E}{\mathbb{E}}
\newcommand{\eps}{\varepsilon}
\newcommand{\dd}{\,\mathrm{d}}
\begin{document}
\section*{UnsolvedMath ID: 10400219}
\textbf{Source problem code:} catalogue entry
\subsection{Definitions and mathematical statement}
For a knot \(K\subset S^3=\partial B^4\), let \(g_s(K)\) be its four-genus: in the smooth convention, the minimum genus of a smooth, connected, compact, oriented surface properly embedded in \(B^4\) with boundary \(K\). A slice knot bounds a properly embedded disk in the same category, equivalently has four-genus zero. A crossing switch interchanges the overpassing and underpassing arcs at one crossing in a knot diagram; ambient isotopies may be used between switches.

Askitas asks whether every knot \(K\) can be changed to a slice knot using \(g_s(K)\) crossing switches. Writing
\[u_s(K)=\min\{n\geq0:\text{a sequence of }n\text{ crossing switches changes }K\text{ to a slice knot}\},\]
the usual ``using at most'' interpretation is
\[\forall K,\quad u_s(K)\leq g_s(K).\]
The customary crossing-change cobordism inequality gives \(g_s(K)\leq u_s(K)\), so in that convention the question is whether \(u_s(K)=g_s(K)\) for all knots. Equivalently, seek a sequence \(K=K_0,\ldots,K_n\) with \(n\leq g_s(K)\), each step a crossing switch up to isotopy, and \(g_s(K_n)=0\). The terminal knot may be any slice knot.

The source does not explicitly state the category for four-genus or slicing. A locally flat topological interpretation would replace the surfaces and disks throughout, and cannot be silently identified with the smooth question. It also does not clarify a literal ``exactly'' versus the customary minimal-number interpretation; the displayed formulation makes that choice visible rather than imposing diagram restrictions absent from the source.
\subsection{Short English statement}
Can every knot be made slice with as many crossing switches as its four-genus? In the usual minimal-number formulation this asks whether the slicing number always equals the four-genus, with a consistent smooth or topological category.
\subsection{Sources}
\begin{itemize}
\item Local UnsolvedMath registry, permanent ID \(10400219\), exact source code \texttt{catalogue entry}.
\item N. Askitas, attribution in Ohtsuki, \emph{Problems on invariants of knots and 3-manifolds} (2002), Problem 12.14, PDF page 166: \url{https://msp.org/gtm/2002/04/gtm-2002-04-024s.pdf}. Registry access 2026-07-29, pdf-font-structure extraction.
\end{itemize}
\end{document}
